\section{Pressure control of moving peaks and the untruncated raw law}
\label{sec:v57-transition-tails}
A bounded central region contains only a finite fraction of the exact
return labels. Passing from its local variation law to the whole return
measure therefore requires tightness of the arithmetic reference itself.
Continuity of a moving spectral drift is not sufficient: its displacement
is multiplied by the square root of the collision count. We obtain the
required bound from a positive exponential moment of the same collision
observable. No estimate at a growing frequency band is used.

Retain the peak data $a_j,z_j,\kappa_j,\mu_j,H_j,B_j$ and the exact
kernel $\mathcal L_{m,R}$ of \eqref{eq:v41-transition-kernel}. Write
$\mu_R=(0,0,\bar\tau_R,c)$, $c=91/(10000\pi)$, and
$X_R=\mathsf f_R-\mu_R$. All constants below are uniform on the
original circular parameter family. The finitely many peak charts and
their partition are unchanged.

\subsection{A real exponential tilt dominates each visible branch}
\begin{lemma}[Uniform centered pressure]
\label{lem:v57-centered-pressure}
For some $t_0,C>0$ there is a real function $\mathscr P_R(t)$,
$|t|<t_0$, such that
\begin{equation}\label{eq:v57-positive-pressure}
 0\le\mathscr P_R(t)\le C|t|^2,\qquad
 \int\exp\left(-t\cdot\sum_{l=0}^{m-1}X_R\circ T_R^l\right)d\nu
                       \le C e^{m\mathscr P_R(t)}\quad(m\ge1).
\end{equation}
It is the logarithm of the positive simple eigenvalue of the centered
occupation twist at frequency $it$ near zero.
\end{lemma}
\begin{proof}
Center the analytic full occupation operator by multiplying it by
$e^{-iz\cdot\mu_R}$. The near-zero splitting and its physical pairing
in Section~\ref{sec:v37-occupation-local}, together with the full
occupation realization of Section~\ref{sec:v40-occupation-powers},
give, on a common small complex ball,
\[
 \int e^{-t\cdot S_mX_R}\,d\nu
       =\lambda_{0,R}(it)^m A_{0,R}(it)+O(\rho^m),\qquad \rho<1.
\]
Here $A_{0,R}(0)=1$, so the amplitude is positive and bounded away
from zero for real small $t$. Conjugation and simplicity make the
eigenvalue real; closeness to one makes it positive and larger than
a fixed number exceeding $\rho$. Jensen's inequality and
$\int X_Rd\nu=0$ show that the left side is at least one for every
$m$. Taking $m$th roots in the splitting proves
$\lambda_{0,R}(it)\ge1$. Put
$\mathscr P_R(t)=\log\lambda_{0,R}(it)$. Centering removes its
linear term, and the uniform analytic second-derivative bound gives
$\mathscr P_R(t)\le C|t|^2$. The displayed splitting, with the
complement absorbed using $\lambda_{0,R}(it)\ge1$, proves the
moment bound. Only a fixed complex neighborhood of zero is involved.
\end{proof}

\begin{theorem}[Drift is controlled by spectral damping]
\label{thm:v57-damped-drift}
There is $C_D<\infty$ such that every retained peak satisfies
\begin{equation}\label{eq:v57-drift-damping}
                  |\mu_j(R)-\mu_R|^2\le C_D\kappa_j(R).
\end{equation}
In particular a peak with $m\kappa_j=O(1)$ has a bounded displacement
$\sqrt m(\mu_j-\mu_R)$. The estimate concerns the existing scalar
branches, including through changes of arithmetic type.
\end{theorem}
\begin{proof}
We first work near an actual resonance on a given peak chart. Its
rank-one peripheral projection is nonzero and has the physical
representation established in Section~\ref{sec:v40-fixed-count-local}.
Choose bounded smooth initial and terminal functions $a,d$ whose
pairing with this projection is nonzero. Such functions exist by the
physical representation and density of smooth functions in $L^2(\nu)$.
These are proof tests on the full collision space, not section weights.
In particular their pairing is not required to have the possibly
vanishing section amplitude $B_j$.

The scalar-continuity proof of Lemma~\ref{lem:v41-scalar-continuity}
applies also to this pair: finite-time pairings are continuous, and
multiplication by the inverse leading eigenvalue makes their limit
uniform on a smaller complex neighborhood. Hence its amplitude
$A_j^{a,d}(R,z)$ is bounded away from zero there. After centering the
branch, write
\[
 \widetilde\lambda_j(R,z)=e^{-iz\cdot\mu_R}\lambda_j(R,z).
\]
Shrink the imaginary neighborhood so that the centered complementary
powers are bounded by $C\rho_1^m$, $\rho_1<1$. The exact pairing,
valid also at complex frequencies for this bounded observable, gives
\begin{align*}
 |A_j^{a,d}(R,z)|\,|\widetilde\lambda_j(R,z)|^m
 &\le\left|\int a(x)e^{iz\cdot S_mX_R(x)}d(T_R^mx)\,d\nu(x)\right|
                       +C\rho_1^m\\
 &\le\|a\|_\infty\|d\|_\infty
                   \int e^{-t\cdot S_mX_R}\,d\nu+C\rho_1^m,
 \qquad z=a_j(R)+it.
\end{align*}
Lemma~\ref{lem:v57-centered-pressure} and $m$th roots imply the
spectral domination
\begin{equation}\label{eq:v57-spectral-pressure}
 \operatorname{Re}\log\widetilde\lambda_j(R,a_j+it)
                                       \le\mathscr P_R(t).
\end{equation}
No order structure on the anisotropic space is presumed: domination
was established on physical scalar pairings with a nonzero amplitude.

At the real maximum the centered logarithmic gradient is
$i v_j$, where $v_j=\mu_j-\mu_R$ is real. The uniform complex
Taylor estimate and \eqref{eq:v57-spectral-pressure} give
\[
             -\kappa_j-v_j\cdot t\le C_1|t|^2
                    \qquad(|t|<t_1).
\]
All drifts are bounded on the compact chart supports. Enlarge $C_1$
so that $|v_j|/(2C_1)<t_1$ and substitute $t=-v_j/(2C_1)$.
This yields $|v_j|^2\le4C_1\kappa_j$.

For completeness, the local tests just chosen need not norm a whole
chart. The zero set of $\kappa_j$ on its compact parameter support is
compact. Cover it by finitely many of the neighborhoods above. On the
remaining compact set $\kappa_j$ has a positive minimum, while $v_j$
is bounded; their quotient therefore satisfies the same conclusion
with a larger constant. If the zero set is empty this last argument
suffices. A maximum over the finite peak family proves
\eqref{eq:v57-drift-damping} without replacing its partition.
\end{proof}

\subsection{A Gaussian envelope centered at the physical mean}
\begin{lemma}[Complex Gaussian and label tails]
\label{lem:v57-transition-envelope}
There are $C,a>0$ such that for every $m\ge1$ and every real target
$y=(k_1,k_2,u,n)$,
\begin{equation}\label{eq:v57-Gaussian-envelope}
 |\mathcal L_{m,R}(y)|+|\mathcal K_{m,R}(y)|
                 \le C\exp\left(-a\frac{|y-m\mu_R|^2}{m}\right).
\end{equation}
For fixed return count $n\ge2$ define
\[
 V_{n,R}(k,m,u)=\frac{(k_1,k_2,m-n/c,u-n\bar\tau_R/c)}{\sqrt n},
 \quad \gamma_{n,R}(k,m,u)=m^{-2}\mathcal L_{m,R}(k_1,k_2,u,n).
\]
With $m\ge n$, $k\in\Z^2$ and $u>0$, one has
\begin{align}
 \sum_{m,k}\int_0^\infty|\gamma_{n,R}|\,du&\le C,
                                      \label{eq:v57-total-reference}\\
 \sum_{m,k}\int_{\{|V_{n,R}|>M,\ u>0\}}|\gamma_{n,R}|\,du
                       &\le C e^{-aM^2}+C e^{-an}
                                      \quad(M>0).
                                      \label{eq:v57-reference-tail}
\end{align}
The constants may change between these displays.
\end{lemma}
\begin{proof}
On the compact matrix family in \eqref{eq:v41-peak-data},
$\operatorname{Re}H\ge a_0I$ and $\|H\|\le A$. The Gaussian integral
in \eqref{eq:v41-complex-gaussian}, by analytic continuation of
completion of the square on $\operatorname{Re}H>0$, equals
$(2\pi)^{-2}(\det H)^{-1/2}\exp(-z^{\mathsf T}H^{-1}z/2)$.
The determinant branch is continued from positive real matrices.
For real $z$, put $w=H^{-1}z$. Then
\[
 \operatorname{Re}(z^{\mathsf T}H^{-1}z)
   =\operatorname{Re}(w^*H^*w)\ge a_0|w|^2
                                      \ge (a_0/A^2)|z|^2.
\]
Also every singular value of $H$ is at least $a_0$. Consequently
$|G_H(z)|\le C e^{-a_1|z|^2}$ uniformly.

Set $Z=(y-m\mu_R)/\sqrt m$ and $d_j=\sqrt m(\mu_j-\mu_R)$.
The modulus of the $j$th transition summand is at most
$C\exp(-m\kappa_j-a_1|Z-d_j|^2)$. By
Theorem~\ref{thm:v57-damped-drift}, $|d_j|^2\le C_Dm\kappa_j$.
The inequality $|Z|^2\le2|Z-d_j|^2+2C_Dm\kappa_j$ proves
\eqref{eq:v57-Gaussian-envelope} after summing the finitely many
bounded amplitudes, including the fixed factor $1/c$.

Summing the two displacement Gaussians and integrating the roof
Gaussian bounds the contribution at count $m$ by
\begin{equation}\label{eq:v57-count-tail-sum}
       C m^{-1/2}\exp\left(-a\frac{(n-cm)^2}{m}\right).
\end{equation}
Indeed $\sum_{k\in\Z^2}e^{-a|k|^2/m}\le Cm$ and the roof
integral over the whole real line is $C\sqrt m$.
For $n/(2c)\le m\le2n/c$, this sum is bounded by
$C n^{-1/2}\sum_m e^{-a'(m-n/c)^2/n}\le C$.
For $n\le m<n/(2c)$ its total is $C\sqrt n e^{-a'n}$, and for
$m>2n/c$ it is bounded by $C\sum_{m>2n/c}e^{-a'm}$.
Absorb the polynomial in an exponential. This proves
\eqref{eq:v57-total-reference}.

On the middle range $m\asymp n$, the linear identities
$m-n/c=-(n-cm)/c$ and
$u-n\bar\tau_R/c=(u-m\bar\tau_R)+\bar\tau_R(m-n/c)$ give
$|V_{n,R}|\le C|Z|$. On $|V_{n,R}|>M$, spend half the Gaussian
exponent to obtain $Ce^{-a'M^2}$, and sum the remaining half as
above. The two outer count ranges have total $Ce^{-a'n}$ whether
or not the target restriction is imposed. This proves
\eqref{eq:v57-reference-tail}. No central cutoff occurs in
$\gamma_{n,R}$ itself.
\end{proof}

\subsection{The complete fixed-return mixed measure}
Let $\mathcal O_n=\Z^2\times\{n,n+1,\ldots\}\times(0,\infty)$,
with counting measure in its first three coordinates and Lebesgue
measure in its last coordinate; denote this product measure by $d\zeta$.
The law $\mathsf P_{n,R}$ of the original $J_{n,R}$ on this space has
density $p_{n,R}(k,m,u)$. It is a probability, not the mass of a
fixed-collision section intersection. Put
\begin{equation}\label{eq:v57-full-reference}
 Z_{n,R}^{\rm ref}=\int_{\mathcal O_n}(\gamma_{n,R})_+\,d\zeta,
 \qquad
 d\mathsf A_{n,R}=\frac{(\gamma_{n,R})_+}{Z_{n,R}^{\rm ref}}d\zeta.
\end{equation}
The second definition is used when its denominator is positive.

\begin{theorem}[Untruncated arithmetic raw law in total variation]
\label{thm:v57-global-raw-TV}
For the same original return index $n$,
\begin{equation}\label{eq:v57-full-L1}
 E_n:=\sup_R\int_{\mathcal O_n}
              |p_{n,R}(k,m,u)-\gamma_{n,R}(k,m,u)|\,d\zeta
                                                   \longrightarrow0.
\end{equation}
Consequently $Z_{n,R}^{\rm ref}\to1$ uniformly in $R$, the reference
in \eqref{eq:v57-full-reference} is eventually defined for every
radius, and
\begin{equation}\label{eq:v57-full-TV}
 \sup_R d_{\rm TV}(\mathsf P_{n,R},\mathsf A_{n,R})\le E_n\to0.
\end{equation}
Here $d_{\rm TV}$ is one half of the variation mass. Every displacement,
collision count and positive roof is included; no central restriction,
source guard, roof smoothing or average over $n$ remains in the result.
The negative transition mass also tends uniformly to zero.
\end{theorem}
\begin{proof}
Let $e_m$ be the unit-window supremum in
Theorem~\ref{thm:v49-raw-local-TV}. It is finite and tends to zero.
Replace it by its decreasing tail supremum when convenient. On
$|V_{n,R}|\le M$, for a fixed $M$, one has $m=n/c+O_M(\sqrt n)$.
There are $O_M(n^{3/2})$ discrete label triples $(k,m)$, and each
roof interval is covered by $O_M(\sqrt n)$ unit intervals. The
local theorem, with the signed kernel still retained, therefore gives
\begin{equation}\label{eq:v57-central-error}
 \int_{\{|V_{n,R}|\le M\}}|p_{n,R}-\gamma_{n,R}|\,d\zeta
                         \le C_M\sup_{m\ge n}e_m\longrightarrow0.
\end{equation}
These sums are taken at a fixed return index. They are the total
variation norm of its joint measure, not an average of different
return laws.

The fourth-moment estimate \eqref{eq:actual-marked-fourth}, with mark
zero, yields
$\int|J_{n,R}-n\bar G_R|^4d\nu_R^*\le Cn^2$.
Thus the original mass of $|V_{n,R}|>M$ is at most $CM^{-4}$.
Lemma~\ref{lem:v57-transition-envelope} supplies the reference tail.
Combining these bounds gives, for each fixed $M$ and all sufficiently
large $n$,
\[
 E_n\le C_M\sup_{m\ge n}e_m+CM^{-4}+Ce^{-aM^2}+Ce^{-an}.
\]
First let $n\to\infty$ and then $M\to\infty$. This proves
\eqref{eq:v57-full-L1}. The preliminary central cutoff was only
in the proof and has disappeared from both measures.

Since $p_{n,R}\ge0$, replacing $\gamma$ by $\gamma_+$ decreases
its pointwise error against $p$. Therefore
$|Z_{n,R}^{\rm ref}-1|\le E_n$, and
\[
 \int\left|p_{n,R}-\frac{\gamma_+}{Z_{n,R}^{\rm ref}}\right|d\zeta
 \le\int|p_{n,R}-\gamma_+|d\zeta+|Z_{n,R}^{\rm ref}-1|
 \le2E_n.
\]
This proves \eqref{eq:v57-full-TV}. On $\gamma<0$ its negative
mass is bounded by $|p-\gamma|$, giving the last assertion.
No pointwise arithmetic floor is required to normalize this whole law.
\end{proof}

\begin{corollary}[The fixed-radius residue is retained globally]
\label{cor:v57-fixed-radius-raw}
At each fixed radius define on $\mathcal O_n$
\[
 q^0_{n,R}(k,m,u)=n^{-2}\mathfrak a_R(k,n,m)
                                      g_{D_R}(V_{n,R}(k,m,u)).
\]
Its mass tends to one, and the probability obtained by dividing it
by its mass is asymptotic in total variation to $\mathsf P_{n,R}$.
The coefficient $\mathfrak a_R$ is the actual finite residue, including
its zero classes; it has not been replaced by one. This specialization
is for fixed $R$, whereas \eqref{eq:v57-full-TV} is radius-uniform.
\end{corollary}
\begin{proof}
On a fixed central region use the fixed-radius specialization in
Theorem~\ref{thm:v41-uniform-transition}. With $L_R$ as in
\eqref{eq:v37-covariance-change}, the centered variables satisfy
$Z=\sqrt{n/m}\,L_RV_{n,R}$ and $m/n\to c^{-1}$ there.
The identity $\Omega_R=cL_RD_RL_R^{\mathsf T}$ and $\det L_R=c$
give
$g_{\Omega_R}(\sqrt cL_Rv)=c^{-3}g_{D_R}(v)$.
Consequently $c m^{-2}g_{\Omega_R}(Z)=n^{-2}g_{D_R}(V)+o(n^{-2})$
uniformly on that region. The finite-radius residue is nonnegative
and bounded. Its Gaussian expression $q^0$ therefore has a Gaussian
tail in $V$ after summing the scaled lattice and integrating the roof,
by the same elementary sums as above. Combine the central comparison
with that tail and the true fourth-moment tail, then normalize exactly
as in Theorem~\ref{thm:v57-global-raw-TV}.
\end{proof}

\paragraph{The positive-height problem.}
The new input is \eqref{eq:v57-drift-damping}; mere continuity at
$\kappa_j=0$ would not give the Gaussian tail used in the proof.
The conclusion integrates the absolute error on the entire exact
mixed record. It does not estimate the essential height of either
physical component in Corollary~\ref{cor:v51-positive-criterion}.
Those two estimates still govern the original unrestricted pointwise
arithmetic local theorem. In particular the whole-law normalization
above does not assign a nonzero probability to a single roof value or
a positive likelihood to a zero arithmetic class.
